Consider, for example, the following two coupling schemes: 1) first $j_1$ and $j_2$ are added to form the total angular momentum $j_{12}$, after which $j_{12}$ and $j_3$ are added to give the final angular momentum $J$; 2) $j_2$ and $j_3$ are added to give $j_{23}$, and then $j_{23}$ is combined with $j_1$ to give $J$. The first scheme corresponds to states in which, besides $j_1,j_2,j_3,J,M$, the quantity $j_{12}$ has a definite value; let their wave functions be denoted by $\psi_{j_{12}JM}$ (the repeated indices $j_1j_2j_3$ being suppressed for brevity). Similarly, denote the wave functions for the second coupling scheme by $\psi_{j_{23}JM}$. In either case the ``intermediate'' angular momentum ($j_{12}$ or $j_{23}$) can, in general, take more than one value. Thus, for fixed $J,M$, there are two distinct sets of states, distinguished respectively by the values of $j_{12}$ and $j_{23}$. By the general rules, the functions in these two sets are related by a certain unitary transformation
